\documentclass[11pt]{article}
\usepackage{amsmath,amssymb,amsthm,hyperref,geometry}\geometry{margin=1in}
\newtheorem{theorem}{Theorem}
\title{A relative parallel-composite criterion}
\date{15 November 2026}\begin{document}\maketitle
\paragraph{Status.} Independently reviewed conditional construction. This supplies a conditional structural route, not a completed feedback/concurrency unifier. The input includes parallel structure on each independent side and compatibility of the crossing; it does not include a parallel compositor for the mixed construction. Deriving the crossing compatibility from geometric hypotheses is a further obligation.

\section{Typed input and direction}
Fix a small monoidal pure category $C$ and Cat-valued endomodules $X:C^{op}\times C\to\mathbf{Cat}$. Use the original relative pseudocoend composition, denoted $XY$, with its invertible pure-middle balances. A parallel operation on $X$ consists of functors
\[
 p_X:X(A,B)\times X(A',B')\longrightarrow X(A\otimes A',B\otimes B'),
\]
pseudonatural for all four pure endpoints, with coherent invertible associativity and unit comparisons and a unit element over $(I,I)$. Native cells remain arrows, not equivalence classes. The relative unit module is $J(A,B)=C(A,B)$ as a discrete category; its parallel product is tensor of pure arrows, with unit $1_I$. For a bicategorical pure base this must instead be its actual pure hom-category, with the corresponding coherent structure; the theorem here fixes the ordinary pure category above. A parallel-colax module map $F:X\to Y$ has natural comparisons
\[
 Fp_X(x,x')\longrightarrow p_Y(Fx,Fx')
\]
and a unit comparison $F(\iota_X)\to\iota_Y$, satisfying the usual associativity/unit equations. These comparison transformations themselves are pseudonatural in the pure endpoints and compatible with all endpoint action constraints. A transformation of such maps is parallel-compatible when its two comparison pastings agree. These requirements are imposed on actual native categories and their endpoint actions.

\paragraph{Parallel structure on a word.}
For two modules with these operations put
\[
 p_{XY}([B;x,y],[B';x',y'])
 =[B\otimes B';p_X(x,x'),p_Y(y,y')].                 \tag{P}
\]
A pure balance at $u:B\to D$ in the first variable maps to the balance at $u\otimes1_{B'}$ followed by the endpoint pseudonaturality comparisons; similarly the other variable uses $1_B\otimes v$. The cross-variable square follows from functoriality of $\otimes$ and the pseudonaturality axioms. Thus the coend universal property, including its modification level, gives an actual functor (P). This is not an assumed commutation of an arbitrary effect with coends. Products of categories preserve the relevant pseudocolimits separately; equivalently the same argument can be made twice by the universal property. Word reassociation and unit equivalences are parallel-compatible because both sides act identically on representatives and balances, up to the specified pure and parallel constraints.

\section{Independent sides}
Let $M,N$ be the original one-sided relative pseudomonads with units $e_M,e_N$ and multiplications $m_M,m_N$. Give each its parallel operation. Require its multiplication to be parallel-colax, with directed comparison
\[
 m_X p_{XX}\Longrightarrow p_X(m_X\times m_X),\qquad X=M,N.       \tag{Q}
\]
Its unit $e_X:J\to X$ is parallel-strong. Require the pseudomonad associativity and unit modifications to be parallel-compatible. Invertibility of (Q) is not required.

Let $t:NM\to MN$ be a relative distributive crossing, with the four unit/multiplication modifications and their coherence equations from the original construction. Give $t$ a parallel-colax comparison
\[
 t p_{NM}\Longrightarrow p_{MN}(t\times t).                    \tag{T}
\]
Require all four distributive modifications to be parallel-compatible. This is an equation involving the independently supplied $M,N,t$ data, not a requirement that an unknown mixed compositor already exist. It is nevertheless a substantive additional hypothesis, not implied by the unstructured crossing theorem.

\begin{theorem}[Conditional composite criterion]
With these inputs, $H=MN$ has parallel operation (P), unit $J\simeq JJ\xrightarrow{e_Me_N}MN$, multiplication
\[
 m_H=(m_Mm_N)(MtN):MNMN\longrightarrow MMNN\longrightarrow MN,
\]
and a derived parallel-colax multiplication. Its pseudomonad unit and associativity modifications are parallel-compatible. Consequently its native realization has a directed sequential/parallel interchanger, natural in every native cell and pure endpoint. The original underlying relative pseudomonad and its native categories are unchanged.
\end{theorem}
\begin{proof}
Whiskering parallel-colax maps by an identity word yields parallel-colax maps: on representatives its comparison is the supplied comparison in the selected slot, with identities in the other slots. Compatibility with coend balance is exactly endpoint pseudonaturality; compatibility across two slots is naturality. Composition of parallel-colax maps has comparison obtained by first applying the outer map to the inner comparison and then using the outer comparison. Its associativity/unit equations follow by substitution of the two respective equations and naturality. Parallel-compatible transformations are closed under vertical and horizontal composition, by pasting their defining equalities.

Apply these facts to $MtN$ and $m_Mm_N$. The resulting mixed comparison is explicitly the crossing comparison (T), whiskered by the two untouched factors, followed by (Q) for $M$ and for $N$. This constructs, rather than assumes, the interchanger. The preceding representative argument proves it respects every native arrow and every pure balance, so it is a transformation of the actual relative modules.

The standard mixed associator is a pasting of the one-sided associators and the two distributive multiplication modifications; the mixed unitors similarly use the one-sided unitors and distributive unit modifications. Each constituent is parallel-compatible by hypothesis. Closure under pasting makes the derived associator and unitors parallel-compatible. Their underlying pentagon and triangles hold by the original distributive-law theorem. Equality of parallel-compatible transformations is equality of their underlying transformations, so these equations remain valid in the structured setting. The associativity and unit equations for the interchanger follow from the parallel-colax axioms just proved. No native comparison or (Q) has been inverted.
\end{proof}

\section{What is and is not obtained}
The interchanger has the direction
\[
 (f\parallel g);(f'\parallel g')\Longrightarrow
 (f;f')\parallel(g;g').
\]
It can be noninvertible even when the crossing itself is invertible. This is stronger than merely having independent juxtaposition, but it does not specify which physical sharing/admission laws a model should implement.

Paquet--Saville's concurrent pseudomonads additionally require certain multiplication comparisons whiskered by pure units to be invertible. Generic parallel-colax structure alone does not imply these. A sufficient separately checkable condition for the present composite is that every factor in the displayed mixed comparison becomes invertible under the relevant pure-unit substitution. Determining a minimal one-sided and crossing condition implying this is left open; the current theorem does not claim that full concurrent-pseudomonad definition automatically follows.

\paragraph{A useful sufficient pure-strand condition.}
Call the side comparison (Q) pure-strand invertible if it is invertible whenever one of its two parallel strands consists of two composable pure arrows, inserted by the side unit. Call (T) pure-strand invertible if it is invertible whenever one parallel argument in $NM$ is a pair of pure arrows inserted by $e_N,e_M$. Assume both side conditions and this crossing condition. Then the mixed interchanger is invertible whenever one whole $H$ strand consists of two composable pure arrows. Indeed substitute their $e_Me_N$ representatives into the explicit three-factor paste. The crossing factor is invertible by its pure-strand condition. The distributive unit comparisons identify the crossed pure factors with the corresponding side units; these comparisons are invertible. Each multiplication factor is therefore an instance of the assumed side pure-strand condition. A composite of these invertible cells is invertible. This is a local sufficient criterion with no mixed interchanger in its assumptions. It is not yet identified with every exact whiskering convention in the published concurrent-pseudomonad definition.

Nor does this theorem add feedback. A feedback operation would need its own compatibility with (Q) and (T), rather than inheriting the symmetric-monoidal state theorem unchanged. The original recovery of native context/content arrows is retained because the underlying $H$ and its pure interface are unchanged; recovery of additional parallel structure in a particular model is a separate comparison.

\paragraph{Literature and priority.} The precedent is composition of structured monads via a structured distributive law. Paquet--Saville, \emph{Effectful semantics in bicategories}, Definition 6.1 and Proposition 6.4, \url{https://philipsaville.co.uk/concurrent-pseudomonads.pdf}, motivates the noninvertible multiplicative comparison and the separate pure-unit condition. Ordinary monoidal-monad results alone do not prove the higher modification-level assertion here. The proposed contribution is the explicit relative/native descent and pasting criterion in the project's presentation; no novelty claim is made.
\end{document}
